\chapter{Fehlersuche}\label{kap:fehlersuche}
\index{Fehlersuche}

Dieser Anhang sammelt Probleme, die in der Praxis vorkommen, mit ihrer Ursache und der
Abhilfe. Die meisten Probleme betreffen nicht das CAD-Programm selbst, sondern die
Python-Umgebung, in der die Plugins rechnen. Eine erste Diagnose liefert fast immer
\menu{Plugins,Settings...} $\to$ \gui{Check environment}.

\section{Installation und Start}

\paragraph{macOS: \enquote{Die App kann nicht geöffnet werden, da der Entwickler nicht
verifiziert werden kann}.}
Die App ist nicht signiert. Öffnen Sie sie beim ersten Mal mit \maus{Rechtsklick} $\to$
\gui{Öffnen} oder erlauben Sie sie unter \gui{Systemeinstellungen} $\to$ \gui{Datenschutz \&
Sicherheit}. Danach startet sie normal (\autoref{kap:installation}).

\paragraph{Windows: SmartScreen warnt vor dem Installationsprogramm.}
Auch das Installationsprogramm ist nicht signiert. Wählen Sie \gui{Weitere Informationen}
$\to$ \gui{Trotzdem ausführen}.

\paragraph{Windows: Menüleiste und Beschriftungen sind kaum lesbar.}
In Versionen vor dem 18.\,11.\,2026 erschien unter Windows mit hellem Systemdesign dunkle
Schrift auf dunklem Grund. Installieren Sie eine aktuelle Version; sie verwendet unabhängig
vom Systemdesign eine dunkle Farbpalette.

\paragraph{Nach dem Start fehlen die Plugins in der Werkzeugspalte.}
\begin{itemize}
  \item Unter \menu{Plugins,Settings...} steht unten, welche Plugins aus welchem Grund nicht
    geladen wurden.
  \item Plugins aus eigenen Ordnern werden nur geladen, wenn sie signiert sind oder der
    Entwicklermodus eingeschaltet ist (\autoref{kap:plugin-verwaltung}). Die mitgelieferten
    Plugins brauchen keine Signatur.
  \item Ein Plugin kann über \menu{Plugins,Show in the side panels} ausgeblendet sein; seine
    Aktionen stehen dann nur im Menü \menu{Plugins}.
\end{itemize}

\section{Python-Umgebung}\index{Python-Umgebung}

\paragraph{\enquote{cannot start Python (\dots)}.}
Der eingetragene Interpreter existiert nicht oder lässt sich nicht starten. Tragen Sie unter
\menu{Plugins,Settings...} einen gültigen Interpreter ein oder legen Sie mit
\gui{Create environment...} eine neue Umgebung an.

\paragraph{Windows: Ein Klick auf \gui{Create environment...} öffnet den Microsoft Store
oder scheitert sofort.}
Unter Windows gibt es Platzhalter \datei{python.exe} im Ordner \datei{WindowsApps}, die nur
den Store öffnen. Die Anwendung überspringt sie, findet dann aber womöglich gar kein Python.
Installieren Sie Python von \url{https://www.python.org/downloads/} und setzen Sie dabei den
Haken \gui{Add python.exe to PATH}. Empfohlen ist Python~3.13, weil es dafür von allen
Rechenpaketen fertige Pakete gibt.

\paragraph{\gui{Check environment} zeigt rote Kreuze.}
Unter der Liste steht der \code{pip}-Befehl, der die fehlenden Pakete installiert. Kopieren Sie
ihn in ein Terminal (Windows: Eingabeaufforderung) und führen Sie ihn aus. Danach erneut
\gui{Check environment}. Eine Zeile wie \enquote{version 0.1.0 installed, >=0.2 needed} heißt,
dass ein Paket zu alt ist; dann hilft \code{pip install -U} mit dem Paketnamen.

\paragraph{\code{pip} versucht zu kompilieren und scheitert.}
Für die gewählte Python-Version gibt es kein fertiges Paket (Wheel). Das passiert vor allem bei
sehr neuen Python-Versionen. Verwenden Sie Python~3.10 bis 3.13. Ohne Compiler installieren
\code{rtt-fusion} und \code{photonpower-fieldtracing} sich aus dem Quellpaket ohne ihre
C++-Teile; sie rechnen dann richtig, aber deutlich langsamer.

\paragraph{\enquote{ModuleNotFoundError: No module named \dots} im Log.}
Ein Paket fehlt in der Umgebung, die das Plugin benutzt. Prüfen Sie, ob unter
\menu{Plugins,Settings...} wirklich die Umgebung eingetragen ist, in die Sie installiert
haben -- es ist leicht, in ein anderes Python zu installieren als das, mit dem die Plugins
laufen. Am sichersten ist der vollständige Pfad des Interpreters:
\begin{lstlisting}
"<Umgebung>\Scripts\python.exe" -m pip install <paket>      (Windows)
<Umgebung>/bin/python -m pip install <paket>                (macOS, Linux)
\end{lstlisting}

\section{Rechnen}

\paragraph{Der Workflow bricht mit einer Meldung über Glaskataloge ab.}
Die Meldung lautet etwa \enquote{the system's glasses come from the glass catalogue(s) \dots,
which are not loaded \dots; add the .agf file under Glass catalogues in the Raytracer panel and
run again}. Das optische System verwendet Gläser aus einem AGF-Katalog, der weder lesbar noch
im Dokument gespeichert ist -- typisch, wenn eine Datei auf einem anderen Rechner geöffnet
wird. Tragen Sie die \datei{.agf}-Datei im Raytracer-Panel unter \gui{Glass catalogues} ein und
starten Sie \gui{Check workflow} oder \gui{Run workflow}. Der Katalog wird dann im Dokument
gespeichert, und die Datei rechnet künftig auch ohne ihn (\autoref{kap:raytracer}).

\paragraph{Ein Anzeigeblock meldet \enquote{dark -- the light focuses into less than a pixel}.}
Das Licht trifft die Kamera in einem Fleck, der kleiner als ein Pixel ist; das Bild bleibt
schwarz. Verkleinern Sie das Bildfeld der Kamera, erhöhen Sie die Auflösung, oder stellen
Sie die Verbindung auf Felder (H) bzw. \code{auto}, damit der Workflow in Fokusnähe mit
Feldern statt Strahlen rechnet (\autoref{kap:workflow}).

\paragraph{Eine Berechnung dauert sehr lange.}
\begin{itemize}
  \item Fortschritt und Zwischenschritte stehen im Log des Plugin-Docks. Mit \keys{Esc}
    brechen Sie ab.
  \item Prüfen Sie, ob die C++-Teile der Rechenpakete installiert sind (Pakete von PyPI für
    Python~3.10--3.13 enthalten sie).
  \item Weniger Strahlen bzw.\ Abtastwerte und eine kleinere Bildauflösung verkürzen die
    Rechnung oft drastisch, ohne das Ergebnis qualitativ zu ändern.
  \item Im Workflow werden unveränderte Blöcke nicht neu berechnet. Ändern Sie nur, was nötig
    ist, und speichern Sie das Dokument: Die Ergebnisse der Anzeigeblöcke werden mitgespeichert
    und beim nächsten Öffnen nicht neu berechnet.
  \item Zeigt eine Anzeige nach einem Update noch das alte Ergebnis (z.\,B.\ ein leeres Bild,
    das eine behobene Version gerechnet hatte), wählen Sie \gui{Recompute} im Kontextmenü des
    Blocks: Er wird mit allem dahinter neu gerechnet.
\end{itemize}

\paragraph{\enquote{plugin '\dots' stopped (exit code \dots)}.}
Der Python-Prozess des Plugins ist beendet worden, zum Beispiel durch einen Fehler in einem
Rechenpaket oder zu wenig Arbeitsspeicher. Das Programm selbst läuft weiter; nur diese Aktion
ist verloren. Die Zeilen davor im Log nennen meist die Ursache. Beim nächsten Aufruf startet
das Plugin neu.

\section{Dateien}

\paragraph{Eine Datei lässt sich nicht öffnen: neuere Version.}
Die Datei wurde mit einer neueren Version des Programms gespeichert, deren Format diese Version
nicht kennt. Aktualisieren Sie das Programm (\menu{Help,Check for updates...}).

\paragraph{STEP-Import: Flächen fehlen oder sind offen.}
Beim Import werden die Flächen zu Körpern vernäht. Liegen Kanten im Ausgangsmodell zu weit
auseinander, bleiben sie offen. Prüfen Sie das Modell im Ausgangsprogramm, oder exportieren Sie
es dort mit höherer Genauigkeit (\autoref{kap:dateien}).

\section{Absturz}\index{Absturzbericht}

Stürzt das Programm ab, schreibt es einen Absturzbericht auf diesen Rechner -- es wird nichts
verschickt. Beim nächsten Start weist es darauf hin. Hängen Sie den Bericht an, wenn Sie das
Problem auf
\url{https://github.com/PhotonPower/Raytatouille-Fusion/issues} melden, und beschreiben Sie,
was Sie gerade getan haben.
